\documentclass[10pt,a4paper]{amsart}
\usepackage[margin=19mm]{geometry}
\usepackage{amsmath,amssymb,mathtools}
\usepackage[scr=rsfs,cal=euler]{mathalfa}
\usepackage{xfrac}
\usepackage[unicode,hidelinks,bookmarks=false]{hyperref}
\newcommand{\ie}{i.e.\ }
\newcommand{\cf}{cf.\ }
\newcommand{\resp}{respectively\ }
\newcommand{\Subsection}{\subsection}
\providecommand{\nobreakdash}{\nobreak\hbox{-}}
\setlength{\emergencystretch}{2em}
\multlinegap0pt
\usepackage[shortlabels]{enumitem}
\usepackage{bm}
\usepackage{amssymb}
\usepackage{mathtools}
\newtheorem{lemma}{Lemma}
\newtheorem{proposition}{Proposition}
\title{The equality between the Erd\H{o}s-Ginzburg-Ziv constant and the short product-one constant for finite nonabelian groups}
\author{Yongke Qu, Guoqing Wang, Yuanlin Li}
\date{Source: arXiv:2608.15568v1 (CC BY 4.0)}
\begin{document}
\maketitle

\noindent\emph{Source and licence.} The equality between the Erd\H{o}s-Ginzburg-Ziv constant and the short product-one constant for finite nonabelian groups. Version arXiv:2608.15568v1 (CC BY 4.0), distributed by arXiv under the Creative Commons Attribution 4.0 International licence, http://creativecommons.org/licenses/by/4.0/, as recorded in the arXiv licence metadata for this exact version and shown on the abstract page https://arxiv.org/abs/2608.15568v1. Full licence notice: \texttt{licenses/arxiv-2608.15568-CC-BY-4.0.txt}.\par\smallskip

\noindent\textit{Editorial scope.} Counted unit: the article's proposition proposition:ifsg=d+exp (source0.tex lines 647--658). The article's complete original proof of it is delivered with it (source0.tex lines 660--720, one \texttt{proof} environment of 61 lines). No mathematics is written, repaired or re-derived: the statement and the proof are the article's own lines, in the article's own notation, and no retained line is substituted or re-flowed.  Retained uncounted as the source-local context the proof needs: lines 503--505.  Nothing was omitted from any retained span, so the package carries no omission record.  No retained line contains a citation, so the wrapper carries no bibliography and no bibliography entry.  No retained line contains a figure environment, an image-inclusion command or drawing code, so no image is delivered and the package declares no asset dependency.  Editorial formatting only, in addition: an AMS \texttt{amsart} preamble that restates the article's own declarations and macros used by the retained lines, namely the environments \texttt{lemma}, \texttt{proposition} on separate counters, together with the standard packages those lines need.  The retained environments print in this document as Lemma 1, Proposition 1 in order of appearance, whereas the article prints the same items with the numbering of its own full document.  The complete native source of the article as distributed in the arXiv e-print for this exact version is bundled unchanged as \texttt{sources/207759/source0.tex}.
\begin{lemma}\label{lemma:universallower}
For every finite group $G$, we have $$s(G)\geq\eta(G)+\exp(G)-1\geq \mathsf d(G)+\exp(G).$$
\end{lemma}

\begin{proposition}\label{proposition:ifsg=d+exp}
Let $G$ be a finite group.  Then $$0\leq E(G)-\bigl(\mathsf d(G)+|G|\bigr)
\leq s(G)-\bigl(\mathsf d(G)+\exp(G)\bigr),$$
and $$s(G)=\mathsf d(G)+\exp(G) \quad \text{if and only if} \quad \ell(G)=1.$$ Moreover, if $\ell(G)=1$ then the following assertions hold.
\begin{enumerate}[label=\rm(\roman*)]
\item For every positive integer $m$,
$s_{m\exp(G)}(G)=\eta(G)+m\exp(G)-1;$
\item $s_{\exp(G)\mathbb N}(G)
=s(G)
=\eta(G)+\exp(G)-1.$
\end{enumerate}
\end{proposition}

\begin{proof}
Let $U\in \mathcal F(G)$ be a product-one free sequence of length $\mathsf d(G)$.

We first prove that, for every
positive integer $m$,
\begin{equation}\label{equation:prescribedbounds}
\mathsf d(G)+m\exp(G)\leq s_{m\exp(G)}(G)\leq s(G)+(m-1)\exp(G).
\end{equation}
 Notice that the
sequence
$U\boldsymbol{\cdot}1^{[m\exp(G)-1]}$
has no product-one subsequence of length $m\exp(G)$. This implies
$s_{m\exp(G)}(G)\geq |U\boldsymbol{\cdot}1^{[m\exp(G)-1]}|+1=\mathsf d(G)+m\exp(G)$.
For the second inequality, let $S$ be a sequence of length
$s(G)+(m-1)\exp(G)$. Applying the definition of $s(G)$ successively, we obtain $m$
mutually disjoint product-one subsequences $T_1,\ldots,T_m$ of $S$, each of length $\exp(G)$.
Then $T_1\boldsymbol{\cdot}\ldots\boldsymbol{\cdot} T_m$ is a product-one subsequence of $S$ of length $m\exp(G)$. This proves \eqref{equation:prescribedbounds}.


Taking
$m=\frac{|G|}{\exp(G)}$ in \eqref{equation:prescribedbounds} and using
$E(G)=s_{m\exp(G)}(G)$, we obtain
$\mathsf d(G)+|G|\leq E(G)\leq s(G)+|G|-\exp(G).$
Subtracting $\mathsf d(G)+|G|$ yields
$$0\leq E(G)-\bigl(\mathsf d(G)+|G|\bigr)
\leq s(G)-(\mathsf d(G)+\exp(G)).$$




We now prove that
$$s(G)=\mathsf d(G)+\exp(G)
\quad \text{if and only if}  \quad
\ell(G)=1.$$  If $\ell(G)=1$, then taking $m=1$ in the definition immediately gives
$s(G)=\mathsf d(G)+\exp(G)$.
Hence, we next assume
$s(G)=\mathsf d(G)+\exp(G)$  and prove $\ell(G)=1$. By \eqref{equation:prescribedbounds}, we derive that for every positive integer $m$, $\mathsf d(G)+m\exp(G)\leq s_{m\exp(G)}(G)\leq s(G)+(m-1)\exp(G)=(\mathsf d(G)+\exp(G))+(m-1)\exp(G)=\mathsf d(G)+m\exp(G)$ and so $s_{m\exp(G)}(G)=\mathsf d(G)+m\exp(G)$. Therefore, $\ell(G)=1$, as required.

Suppose now that $\ell(G)=1$, i.e.,
$s(G)=\mathsf d(G)+\exp(G)$. By
Lemma \ref{lemma:universallower}, we have
\begin{equation}\label{equation:eta=d+1}
\eta(G)=\mathsf d(G)+1
\qquad\text{and}\qquad
s(G)=\eta(G)+\exp(G)-1.
\end{equation}
Combined with the definition of $\ell(G)$,  we have that $s_{m\exp(G)}(G)=\mathsf d(G)+m \exp(G)=\eta(G)+m \exp(G)-1$ for every $m\geq 1$,
which proves \rm(i).

By definition,
$s_{\exp(G)\mathbb N}(G)\leq s(G)$.
On the other hand, the sequence $U\boldsymbol{\cdot}1^{[\exp(G)-1]}$
contains no product-one subsequence whose length is a positive
multiple of $\exp(G)$. Consequently,
$s_{\exp(G)\mathbb N}(G)\geq |U\boldsymbol{\cdot}1^{[\exp(G)-1]}|+1=\mathsf d(G)+\exp(G)=s(G).$
Together with \eqref{equation:eta=d+1}, this gives
$s_{\exp(G)\mathbb N}(G)
=s(G)
=\eta(G)+\exp(G)-1,$
and proves \rm(ii).
\end{proof}

\end{document}
